%==============================================================================
%  Lecture Notes: The Stable Marriage Problem
%  Advanced Algorithms 2026 -- Lecture 1
%
%  These notes accompany the slide deck "The Stable Marriage Problem"
%  (Algorithms and Networks, Johan M. M. van Rooij & Jesper Nederlof).
%
%  ---------------------------------------------------------------------------
%  FIGURES
%  ---------------------------------------------------------------------------
%  This file contains NO images.  Every place where a picture would help is
%  marked with a \figplaceholder{<height>}{<description of the picture>}
%  inside a normal figure environment, so the document compiles as-is and
%  prints a framed grey box describing what should be drawn there.
%
%  To insert a real figure later, simply replace
%      \figplaceholder{45mm}{...description...}
%  by
%      \includegraphics[width=0.8\textwidth]{myfigure}
%  and keep the surrounding \caption / \label untouched, so that all
%  cross-references in the text keep working.
%
%  Search for "figplaceholder" to find them all; there are 17.
%==============================================================================

\documentclass[11pt,a4paper]{article}

\usepackage[T1]{fontenc}
\usepackage[utf8]{inputenc}
\usepackage[margin=3cm]{geometry}
\usepackage{amsmath,amssymb,amsthm}
\usepackage{enumitem}
\usepackage{graphicx}
\usepackage{xcolor}
\usepackage{booktabs}
\usepackage{algorithm}
\usepackage{algpseudocode}
\usepackage[colorlinks=true,linkcolor=blue!50!black,citecolor=blue!50!black,
            urlcolor=blue!50!black]{hyperref}

%------------------------------------------------------------------ theorems --
\theoremstyle{plain}
\newtheorem{theorem}{Theorem}[section]
\newtheorem{lemma}[theorem]{Lemma}
\newtheorem{proposition}[theorem]{Proposition}
\newtheorem{corollary}[theorem]{Corollary}
\newtheorem{observation}[theorem]{Observation}

\theoremstyle{definition}
\newtheorem{definition}[theorem]{Definition}
\newtheorem{example}[theorem]{Example}
\newtheorem{exercise}{Exercise}

\theoremstyle{remark}
\newtheorem{remark}[theorem]{Remark}
\newtheorem*{intuition}{Intuition}

%------------------------------------------------------------------- macros ---
\newcommand{\Men}{\mathcal{M}}
\newcommand{\Women}{\mathcal{W}}
\newcommand{\pr}[1]{\succ_{#1}}          % w \pr{m} w'  reads "m prefers w to w'"
\newcommand{\GS}{\textsc{Gale--Shapley}}
\newcommand{\bigoh}[1]{\mathcal{O}\!\left(#1\right)}

% Figure placeholder: #1 = height, #2 = description of the intended picture.
\newcommand{\figplaceholder}[2]{%
  \setlength{\fboxsep}{6pt}%
  \fcolorbox{gray!60}{gray!8}{%
    \parbox[c][#1][c]{0.86\textwidth}{%
      \centering\small\itshape\color{black!65}%
      [\,figure placeholder\,]\par\vspace{0.6em}#2}}}

\title{\bfseries The Stable Marriage Problem\\[0.4em]
       \large Lecture Notes -- Advanced Algorithms 2026, Lecture 1}
\author{}
\date{}

\begin{document}
\maketitle
\vspace{-3.5em}

\begin{center}
\begin{minipage}{0.86\textwidth}
\small
These notes follow the lecture slides \emph{The Stable Marriage Problem}
by Johan M.\,M.\ van Rooij and Jesper Nederlof (\emph{Algorithms and
Networks}). They cover the same material, but with the arguments written out
in full, together with the informal reasoning that usually gets spoken rather
than projected. Nothing here goes beyond the content of the lecture; the short
outlook in Section~\ref{sec:outlook} is the only part that is explicitly
non-examinable.
\end{minipage}
\end{center}
\vspace{1em}

\tableofcontents
\bigskip

%==============================================================================
\section{Introduction: markets that clear themselves}
%==============================================================================

Most of the algorithmic problems you have met so far come with a built-in
notion of what a good solution is: shortest, cheapest, largest, fastest. You
write down an objective function, and then you optimise it. This lecture is
about a problem where the interesting question is not \emph{how good} a
solution is, but whether it is \emph{self-enforcing} --- whether the people
involved will actually go along with it once you hand it to them.

The setting is deceptively simple. There are two groups of participants of
equal size. Everybody in the first group ranks everybody in the second group,
and vice versa. We must pair them up. There is no global score to maximise;
there is not even an obvious candidate for one, because the participants'
preferences are in direct conflict with each other. What we can ask for is
much weaker, and much more interesting:

\begin{quote}
Is there a pairing that nobody has an incentive to walk away from?
\end{quote}

A pairing with this property is called \emph{stable}. The question of whether
one always exists, and how to find it, was answered in 1962 by David Gale and
Lloyd Shapley in a paper with the memorable title \emph{``College Admissions
and the Stability of Marriage''} \cite{GaleShapley1962}. Their answer is a
short, elegant algorithm --- the \emph{deferred acceptance} or
\emph{Gale--Shapley} algorithm --- which we will develop, analyse and dissect
over the course of this lecture.

\subsection*{Why anyone cares}

The problem is not a puzzle invented for a textbook. It describes a real and
recurring situation: a market in which money does not do the allocating,
because both sides have to agree.

The canonical example, and the one that gave the field its practical urgency,
is the assignment of graduating medical students to hospital residency
positions in the United States. In the 1940s this market had collapsed into
chaos. Hospitals, competing for the best students, kept moving their offers
earlier and earlier --- eventually to two full years before graduation, at
which point neither side had any real information about the other. Offers came
with exploding deadlines of a few hours. Students accepted positions and then
reneged when something better appeared; hospitals did the same. The market was
\emph{unravelling}.

In 1952 the profession adopted a centralised clearinghouse. Everybody submits
a ranked list, a computation is performed, and the resulting assignment is
binding. The crucial design question was: what computation? If the output can
be undermined --- if there is a student and a hospital who would both rather
have each other than what the clearinghouse gave them --- then they will
simply arrange that privately, other participants will notice, and the
clearinghouse loses its authority. Stability is exactly the property that
prevents this. Remarkably, the algorithm the clearinghouse adopted in 1952 was
essentially deferred acceptance, a decade before Gale and Shapley published
it and proved that it works; this was only recognised much later by Alvin Roth
\cite{Roth1984}.

\begin{figure}[ht]
\centering
\figplaceholder{40mm}{A timeline / context figure for the residency match:
``1940s: offers creep to 2 years before graduation, exploding deadlines,
market unravels'' $\to$ ``1952: centralised clearinghouse (NRMP)'' $\to$
``1962: Gale \& Shapley explain why it works'' $\to$ ``2012: Nobel Prize''.
Could be a simple horizontal arrow with four labelled milestones.}
\caption{The stable marriage problem was solved in practice before it was
solved in theory.}
\label{fig:history}
\end{figure}

In 2012 Lloyd Shapley shared the Nobel Memorial Prize in Economic Sciences
with Alvin Roth, ``for the theory of stable allocations and the practice of
market design''. Shapley's citation covers a number of contributions --- he is
at least as famous for the Shapley value in cooperative game theory --- but the
Gale--Shapley algorithm is prominent among them. Roth's share of the prize was
for taking the theory into the field: redesigning the residency match, and
later the New York and Boston public school choice systems and kidney exchange
programmes. David Gale had died in 2008, and the prize is not awarded
posthumously.

\begin{figure}[ht]
\centering
\figplaceholder{45mm}{Portrait of Lloyd Shapley (Nobel Prize in Economics,
2012). Slide~2 of the deck uses the Wikimedia Commons photograph by Bengt
Nyman, \url{http://en.wikipedia.org/wiki/File:Lloyd_Shapley_2_2012.jpg}
--- check the CC licence terms and reproduce the attribution.}
\caption{Lloyd Shapley (1923--2016).}
\label{fig:shapley}
\end{figure}

\subsection*{What this lecture will do}

\begin{itemize}[nosep]
  \item Section~\ref{sec:model} sets up the model and defines stability
        precisely.
  \item Section~\ref{sec:naive} shows why the obvious approach --- repeatedly
        repair whatever is broken --- can run forever, and why the mere
        \emph{existence} of a stable matching is a genuine theorem rather than
        an observation.
  \item Section~\ref{sec:gs} presents the Gale--Shapley algorithm.
  \item Section~\ref{sec:analysis} proves that it terminates within $n^2$
        proposals, that its output is a perfect matching, and that this
        matching is stable. Together these give the existence theorem.
  \item Section~\ref{sec:optimality} asks who benefits. The answer is sharp
        and slightly disturbing: the proposing side gets, simultaneously, the
        best partner it could possibly have in \emph{any} stable matching, and
        the other side gets the worst.
\end{itemize}

\paragraph{A note on terminology.}
The classical language of ``men'' and ``women'' is historical, dating from
Gale and Shapley's original framing, and we keep it because the slides,
the literature and the name of the problem all use it. Mathematically nothing
depends on it: the content is that there are two disjoint sides, and that
matches are only ever formed between the sides and never within one. Read
``students and hospitals'' throughout if you prefer.

%==============================================================================
\section{The model}\label{sec:model}
%==============================================================================

\subsection{Definitions}

\begin{definition}[Instance]
An instance of the \emph{stable marriage problem} consists of a set
$\Men = \{m_1,\dots,m_n\}$ of $n$ men and a set
$\Women = \{w_1,\dots,w_n\}$ of $n$ women, together with, for every
participant, a \emph{preference list}: a strict total order on all $n$
members of the opposite side.
\end{definition}

We write $w \pr{m} w'$ to mean ``man $m$ strictly prefers $w$ to $w'$'',
i.e.\ $w$ appears before $w'$ on $m$'s list, and symmetrically $m \pr{w} m'$
for the women. Three features of this definition deserve to be pointed out,
because all three will matter later:

\begin{itemize}
  \item \textbf{The lists are complete.} Everybody ranks \emph{everybody} on
        the other side. Nobody is unacceptable. This is what will let us prove
        that every participant ends up matched.
  \item \textbf{The lists are strict.} There are no ties, no indifference.
        A woman confronted with two men always has a preference. Allowing ties
        makes several of the questions in this lecture considerably harder
        (see Section~\ref{sec:outlook}).
  \item \textbf{The two sides are equally large.} With $|\Men| \neq |\Women|$
        some participants must remain unmatched, and the definitions below need
        adjusting.
\end{itemize}

\begin{definition}[Matching]
A \emph{matching} is a set $\mu$ of $n$ disjoint man--woman pairs; that is, a
perfect matching in the complete bipartite graph $K_{n,n}$ with sides $\Men$
and $\Women$. We write $\mu(m)$ for the partner of $m$ and $\mu(w)$ for the
partner of $w$, so $\mu(m) = w \iff \mu(w) = m$.
\end{definition}

There are $n!$ matchings, so brute force is out of the question for anything
but tiny instances. Now the key definition.

\begin{definition}[Blocking pair, stability]\label{def:stable}
Let $\mu$ be a matching. A pair $(m,w) \notin \mu$ is a \emph{blocking pair}
for $\mu$ if
\[
  w \pr{m} \mu(m)
  \qquad\text{and}\qquad
  m \pr{w} \mu(w),
\]
that is, if $m$ and $w$ each strictly prefer the other to their assigned
partner. A matching is \emph{stable} if it admits no blocking pair.
\end{definition}

\begin{figure}[ht]
\centering
\figplaceholder{42mm}{The anatomy of a blocking pair. Draw two matched pairs
as horizontal solid edges: $m$--$w'$ on top, $m'$--$w$ below. Then draw the
diagonal $m$--$w$ as a dashed/red edge. Annotate the diagonal with two arrows:
``$m$: $w \succ w'$'' and ``$w$: $m \succ m'$''. Caption idea: both endpoints
of the dashed edge want to switch, so they will.}
\caption{A blocking pair: an unmatched couple who both prefer each other to
what the matching gave them.}
\label{fig:blockingpair}
\end{figure}

\begin{intuition}
A blocking pair is a \emph{private deviation that is guaranteed to succeed}.
It takes exactly two people, both of whom benefit, and neither of whom needs
anybody's permission. That is why stability is the right notion of ``a
solution the market will accept'': it is not a statement about how happy people
are, but about whether the assignment can be undermined by a coalition of the
smallest interesting size.

Note also how weak the requirement is. We do not ask that everyone be happy;
we do not ask that the total or average satisfaction be high; we do not even
ask for fairness. We ask only that no \emph{pair} can improve on the outcome
by themselves. This weakness is precisely what makes the existence theorem
possible --- and, as we shall see in Section~\ref{sec:optimality}, stable
matchings can be spectacularly unfair.
\end{intuition}

\begin{remark}[Why single-person deviations are not an issue]
One could imagine also forbidding a single participant from wanting to walk
away. Because our preference lists are complete, nobody is ever worse off
matched than unmatched, so individual deviations never occur and stability
against pairs is the whole story. In models with incomplete lists (where you
may declare some partners unacceptable) one also requires
\emph{individual rationality}, and blocking pairs are defined only over
mutually acceptable couples.
\end{remark}

\subsection{A worked example}\label{sec:example}

\begin{example}\label{ex:main}
Take $n=3$, with $\Men = \{\text{Arie},\text{Bert},\text{Carl}\}$ and
$\Women = \{\text{Ann},\text{Betty},\text{Cindy}\}$, and preference lists as
follows (best first):

\begin{center}
\begin{tabular}{@{}lccc@{}}
\toprule
\textbf{Men} & 1st & 2nd & 3rd \\
\midrule
Arie  & Ann   & Cindy & Betty \\
Bert  & Betty & Ann   & Cindy \\
Carl  & Betty & Ann   & Cindy \\
\bottomrule
\end{tabular}
\hspace{2em}
\begin{tabular}{@{}lccc@{}}
\toprule
\textbf{Women} & 1st & 2nd & 3rd \\
\midrule
Ann   & Bert & Carl & Arie \\
Betty & Carl & Bert & Arie \\
Cindy & Bert & Carl & Arie \\
\bottomrule
\end{tabular}
\end{center}
\end{example}

Poor Arie: every woman ranks him last. This is a good instance to keep in mind,
because it makes vivid that stability has nothing to do with fairness. Arie is
universally unwanted and no algorithm can fix that; the most we can hope for
is that he ends up with someone who cannot do better than him \emph{given
what everybody else is doing}.

\begin{figure}[ht]
\centering
\figplaceholder{55mm}{The instance of Example~\ref{ex:main} drawn as a
bipartite graph. Men in a column on the left (Arie, Bert, Carl), women on the
right (Ann, Betty, Cindy). Each man's preference list written next to his node
in order, likewise for the women. Optionally draw all nine potential edges
very faintly to emphasise that any of the $3! = 6$ perfect matchings is a
candidate.}
\caption{The instance of Example~\ref{ex:main}.}
\label{fig:instance}
\end{figure}

Consider first the matching
\[
  \mu_{\mathrm{bad}} =
  \{(\text{Arie},\text{Cindy}),\ (\text{Bert},\text{Betty}),\
    (\text{Carl},\text{Ann})\}.
\]
This is \emph{not} stable. Look at Carl and Betty. Carl is matched to Ann, but
Betty is his first choice, so $\text{Betty} \pr{\text{Carl}} \text{Ann}$.
Betty is matched to Bert, but Carl is \emph{her} first choice, so
$\text{Carl} \pr{\text{Betty}} \text{Bert}$. They both prefer each other to
their current partners: $(\text{Carl},\text{Betty})$ is a blocking pair, and
in any real market these two would simply pair up and leave the mechanism
looking foolish.

Now consider
\[
  \mu^{*} =
  \{(\text{Arie},\text{Cindy}),\ (\text{Bert},\text{Ann}),\
    (\text{Carl},\text{Betty})\}.
\]
This one \emph{is} stable. To check this by hand, it suffices to examine, for
each man, only the women he strictly prefers to his current partner --- if he
is content, he cannot be half of a blocking pair.

\begin{itemize}[nosep]
  \item \textbf{Carl} has Betty, his first choice. Content.
  \item \textbf{Bert} has Ann, his second choice; he would rather have Betty.
        But Betty has Carl, whom she ranks first, so she will not entertain
        Bert. Not blocking.
  \item \textbf{Arie} has Cindy, his second choice; he would rather have Ann.
        But Ann has Bert, whom she ranks first. Not blocking.
\end{itemize}

No blocking pair exists, so $\mu^*$ is stable. Notice how the verification
worked: every time a man wanted to trade up, the woman in question was already
holding somebody she liked better. That tension --- men reaching downward
through their lists while women climb upward through theirs --- is the engine
of everything that follows.

\begin{figure}[ht]
\centering
\figplaceholder{50mm}{Side-by-side comparison of $\mu_{\mathrm{bad}}$ (left)
and $\mu^*$ (right) on the bipartite graph. On the left, highlight the
blocking edge Carl--Betty in red with the two ``I prefer you'' arrows; on the
right, no red edge. Same node layout in both panels so the eye can compare.}
\caption{An unstable and a stable matching on the same instance. The red edge
on the left is a blocking pair.}
\label{fig:stable-vs-unstable}
\end{figure}

\begin{remark}[Checking stability is easy; finding a stable matching is the
problem]
Given a matching, we can verify stability in $\bigoh{n^2}$ time by checking all
pairs (with a little care, in $\bigoh{n^2}$ total even naively, since there
are only $n^2$ pairs and each check is $\bigoh{1}$ given a precomputed rank
table). So the problem is squarely in the ``easy to check'' regime. What is
not obvious is how to \emph{find} such a matching --- or whether one exists at
all.
\end{remark}

%==============================================================================
\section{Two false starts}\label{sec:naive}
%==============================================================================

Before presenting the algorithm that works, it is worth spending time on the
approach that does not. This is not just for historical colour: understanding
precisely \emph{why} the naive method fails tells us what the successful
algorithm has to get right.

\subsection{The soap-opera algorithm}

Here is the most natural idea imaginable. A matching is bad exactly when it
has a blocking pair. So: find a blocking pair, and fix it.

\begin{algorithm}[ht]
\caption{Greedy repair (the ``soap series'' algorithm)}
\label{alg:greedy}
\begin{algorithmic}[1]
  \State start from an arbitrary matching $\mu$
  \While{$\mu$ has a blocking pair $(m,w)$}
    \State let $m' = \mu(w)$ and $w' = \mu(m)$
    \State replace the pairs $(m,w')$ and $(m',w)$ by $(m,w)$ and $(m',w')$
  \EndWhile
  \State \Return $\mu$
\end{algorithmic}
\end{algorithm}

The blocking couple gets together; their two jilted ex-partners are left over,
and --- this being a soap series --- console themselves with each other. Every
step makes two people strictly happier. Surely this must get somewhere?

It need not. The two people who were made happy are not the only ones affected:
the two who were dumped are typically made \emph{much} worse off, and they will
generate new blocking pairs of their own. There is no potential function that
obviously decreases, and indeed the process can cycle forever.

\begin{example}[A cycling instance]\label{ex:cycle}
Take three men $1,2,3$ and three women $A,B,C$ with the following preferences
(best first):
\begin{center}
\begin{tabular}{@{}lccc@{}}
\toprule
\textbf{Men} & 1st & 2nd & 3rd \\
\midrule
$1$ & $A$ & $C$ & $B$ \\
$2$ & $C$ & $A$ & $B$ \\
$3$ & $A$ & $B$ & $C$ \\
\bottomrule
\end{tabular}
\hspace{2em}
\begin{tabular}{@{}lccc@{}}
\toprule
\textbf{Women} & 1st & 2nd & 3rd \\
\midrule
$A$ & $2$ & $1$ & $3$ \\
$B$ & $3$ & $1$ & $2$ \\
$C$ & $1$ & $2$ & $3$ \\
\bottomrule
\end{tabular}
\end{center}
Starting from $\{1A,\,2B,\,3C\}$, resolve blocking pairs in this order:
\begin{gather*}
\{1A,2B,3C\}
\;\xrightarrow{\ (2,A)\ }\;
\{1B,2A,3C\}
\;\xrightarrow{\ (2,C)\ }\;
\{1B,2C,3A\}\\[0.3em]
\;\xrightarrow{\ (1,C)\ }\;
\{1C,2B,3A\}
\;\xrightarrow{\ (1,A)\ }\;
\{1A,2B,3C\}.
\end{gather*}
We are back where we started, and the process repeats forever.
\end{example}

Let us verify the first step, and leave the rest to the reader (this is worth
doing once by hand). In $\{1A,2B,3C\}$, is $(2,A)$ a blocking pair? Man $2$ is
matched to $B$, his last choice, and prefers $A$; woman $A$ is matched to $1$,
and ranks $2$ above $1$. Yes. Resolving it marries $2$ to $A$, leaving $1$ and
$B$ --- the two abandoned partners --- to each other. That gives
$\{1B, 2A, 3C\}$.

\begin{figure}[ht]
\centering
\figplaceholder{60mm}{The cycle of Example~\ref{ex:cycle} drawn as a directed
4-cycle. Place the four matchings $\{1A,2B,3C\}$, $\{1B,2A,3C\}$,
$\{1B,2C,3A\}$, $\{1C,2B,3A\}$ at the corners of a square, each drawn as a
small bipartite picture. Label each arrow with the blocking pair that is
resolved: $(2,A)$, $(2,C)$, $(1,C)$, $(1,A)$. Big circular arrow in the middle
to emphasise that this loops forever.}
\caption{Greedy repair can cycle: resolving blocking pairs one at a time
returns to its starting point.}
\label{fig:cycle}
\end{figure}

\begin{intuition}
The lesson is about \emph{commitment}. In the greedy algorithm every match is
provisional in a destructive way: making one couple happy actively injures
another couple, and the damage propagates. What we need instead is a process
with a built-in ratchet --- some quantity that moves in one direction only and
that therefore cannot return to a previous state. The Gale--Shapley algorithm
supplies exactly that, in the form of an asymmetry between the two sides: one
side only ever descends its preference list, the other only ever ascends.
\end{intuition}

\begin{remark}[The greedy algorithm is not hopeless, only unreliable]
This example shows that \emph{some} sequences of repairs cycle. It does not
show that all do. In fact Roth and Vande Vate proved that from any matching
there always \emph{exists} a finite sequence of blocking-pair resolutions
leading to a stable matching, so that if at each step one picks a blocking
pair at random, the process reaches stability with probability~$1$. That is a
beautiful result, but it is not an algorithm with any useful running-time
guarantee, and it does not help us prove that a stable matching exists in the
first place --- one has to know that already. Off-syllabus; mentioned only so
that you do not draw a stronger conclusion from Example~\ref{ex:cycle} than it
supports.
\end{remark}

\subsection{Existence is not free}\label{sec:roommates}

Here is a trap worth walking into deliberately. After staring at small
examples for a while, one starts to feel that a stable matching must surely
always exist --- how could a market fail to have \emph{any} outcome that
nobody can undermine? But this feeling is wrong in general. It is true for the
two-sided problem, and that is a theorem with content.

To see that the theorem has content, drop the two-sidedness. In the
\emph{stable roommates problem} there are $2n$ people, each of whom ranks all
the others, and we pair them up without any bipartite structure. Everything
else --- blocking pairs, stability --- is defined exactly as before.

\begin{example}[No stable matching]\label{ex:roommates}
Four people $a,b,c,d$ with preferences
\[
  a:\ b \succ c \succ d, \qquad
  b:\ c \succ a \succ d, \qquad
  c:\ a \succ b \succ d, \qquad
  d:\ \text{anything}.
\]
The three people $a,b,c$ form a rock--paper--scissors cycle: $a$ prefers $b$,
$b$ prefers $c$, $c$ prefers $a$; and all three rank the unfortunate $d$ last.

Any matching pairs $d$ with one of $a,b,c$ and the other two with each other.
Say $d$ is matched to $a$, so $b$ and $c$ are matched. Now $a$ is matched to
his least-favourite person, so $a$ will block with \emph{anybody} who is
willing --- and $c$ is willing, because $c$ ranks $a$ first and is currently
matched to $b$. So $(a,c)$ blocks. By the cyclic symmetry of the instance the
same argument works whichever of $a,b,c$ is stuck with $d$. Hence no stable
matching exists.
\end{example}

\begin{figure}[ht]
\centering
\figplaceholder{45mm}{The roommates counterexample. Draw $a$, $b$, $c$ as a
triangle with cyclic arrows $a \to b \to c \to a$ labelled ``prefers'', and
$d$ off to the side, with all three arrows pointing at $d$ labelled ``last
choice''. Beside it, show one of the three matchings (say $\{ad, bc\}$) with
the blocking edge $ac$ highlighted.}
\caption{The stable roommates problem: an instance with no stable matching at
all.}
\label{fig:roommates}
\end{figure}

So existence genuinely uses the bipartite structure of the marriage problem.
Keep Example~\ref{ex:roommates} in mind as a control: any proof of existence
for stable marriage that would also work for roommates must be wrong.

%==============================================================================
\section{The Gale--Shapley algorithm}\label{sec:gs}
%==============================================================================

\begin{theorem}[Gale--Shapley, 1962]\label{thm:main}
There is an algorithm which, given a stable marriage instance on $n{+}n$
participants, outputs a stable matching using at most $n^2$ proposals. In
particular, every instance admits a stable matching.
\end{theorem}

The second sentence is worth pausing on. It is a purely mathematical
statement --- no algorithms, no computation --- and yet the only proof we
shall give of it is by exhibiting an algorithm and running it. This is a
recurring and rather satisfying phenomenon: an efficient algorithm is often
the cleanest available proof of an existence theorem.

\subsection{The algorithm}

The idea is \emph{deferred acceptance}. Men propose, in order of their own
preference. Women never accept anything permanently: they hold on to the best
offer they have received so far and remain free to trade up later. An
engagement can always be broken by the woman, never by the man.

\begin{algorithm}[ht]
\caption{Gale--Shapley (men proposing)}
\label{alg:gs}
\begin{algorithmic}[1]
  \State fix an arbitrary ordering of the men; initially everybody is free
  \While{some man is unmatched}
    \State let $m$ be the first unmatched man in the ordering
    \State let $w$ be the highest woman on $m$'s list to whom $m$ has not
           yet proposed
    \State $m$ proposes to $w$
    \If{$w$ is free}
      \State $w$ and $m$ become engaged
    \ElsIf{$w$ prefers $m$ to her current partner $m''$}
      \State $w$ and $m$ become engaged; $m''$ becomes free
      \Comment{$m''$ is \emph{jilted}}
    \Else
      \State $w$ rejects $m$; $m$ stays free
    \EndIf
  \EndWhile
  \State \Return the set of engaged pairs
\end{algorithmic}
\end{algorithm}

Three points of clarification, each of which is a common source of confusion:

\begin{itemize}
  \item \textbf{Rejection is permanent, acceptance is not.} Once $w$ has
        turned $m$ down, $m$ will never propose to her again --- he moves on
        down his list. But an engagement is only ever provisional; $w$ may
        break it the moment somebody better comes along. This asymmetry is the
        whole trick.
  \item \textbf{``Better off with $m$ than with her current partner''} means
        strictly better according to \emph{her} preference list. She does not
        consider how the men feel about it, and she does not consider anybody
        outside the current proposal.
  \item \textbf{The ordering of the men is arbitrary.} The algorithm as stated
        picks the first free man in a fixed order, but any free man will do,
        and we will prove in Section~\ref{sec:optimality} that the output does
        not depend on the choice. This is a strong and surprising statement:
        a completely non-deterministic procedure with a canonical answer.
\end{itemize}

\begin{figure}[ht]
\centering
\figplaceholder{48mm}{One iteration of the algorithm, in three small panels.
Panel 1: free man $m$ with his list, arrow pointing at the next unproposed
woman $w$; $w$ currently engaged to $m''$. Panel 2: $w$ compares $m$ and $m''$
on her list. Panel 3(a): $m \succ_w m''$, so the edge to $m''$ is cut (dashed,
$m''$ drops back into the pool of free men) and a new edge to $m$ is drawn.
Panel 3(b): otherwise $m$ is rejected and returns to the pool with his pointer
advanced by one.}
\caption{A single proposal step. The woman always ends up at least as well off
as before; the man's pointer only ever moves down his list.}
\label{fig:step}
\end{figure}

\subsection{Running the algorithm}\label{sec:trace}

Let us run Algorithm~\ref{alg:gs} on Example~\ref{ex:main}, with the men
ordered Arie, Bert, Carl.

\begin{center}
\small
\begin{tabular}{@{}clll@{}}
\toprule
\# & Proposal & Outcome & Engagements afterwards \\
\midrule
1 & Arie $\to$ Ann   & free; accepts
  & (Arie,\,Ann) \\
2 & Bert $\to$ Betty & free; accepts
  & (Arie,\,Ann), (Bert,\,Betty) \\
3 & Carl $\to$ Betty & prefers Carl; \textbf{Bert jilted}
  & (Arie,\,Ann), (Carl,\,Betty) \\
4 & Bert $\to$ Ann   & prefers Bert; \textbf{Arie jilted}
  & (Bert,\,Ann), (Carl,\,Betty) \\
5 & Arie $\to$ Cindy & free; accepts
  & (Arie,\,Cindy), (Bert,\,Ann), (Carl,\,Betty) \\
\bottomrule
\end{tabular}
\end{center}

After five proposals everybody is matched and the algorithm stops, returning
\[
  \mu_{\Men} =
  \{(\text{Arie},\text{Cindy}),\ (\text{Bert},\text{Ann}),\
    (\text{Carl},\text{Betty})\} = \mu^{*},
\]
the stable matching we verified by hand in Section~\ref{sec:example}.

Trace the dynamics once more and watch the two ratchets turn. Bert starts at
the top of his list (Betty), is displaced, and moves down to Ann. Arie starts
at the top of his list (Ann), is displaced, and moves down to Cindy. Meanwhile
Betty goes from Bert to Carl --- upward on her list --- and Ann goes from Arie
to Bert, also upward. \emph{Every single event in the execution moves some man
down and some woman up.} Nothing ever moves the other way. That is the
property the soap-opera algorithm lacked, and it is what we will now turn into
proofs.

\begin{figure}[ht]
\centering
\figplaceholder{65mm}{The five-step trace as a strip of five small bipartite
snapshots, left to right. In each snapshot: solid edges for current
engagements, a bold arrow for the proposal being made, and a red cross on the
edge that gets broken. Below each snapshot write the proposal, e.g.\
``Carl $\to$ Betty: accepted, Bert jilted''.}
\caption{Execution of Gale--Shapley on Example~\ref{ex:main}.}
\label{fig:trace}
\end{figure}

\begin{remark}[Try a different order]
Run the same instance with the men ordered Carl, Bert, Arie, or let Bert
propose first. The intermediate steps differ, but the final matching is
identical. Theorem~\ref{thm:manoptimal} explains why this is no coincidence.
\end{remark}

%==============================================================================
\section{Analysis}\label{sec:analysis}
%==============================================================================

We now prove Theorem~\ref{thm:main}. There are three things to establish:
the algorithm terminates (and quickly); its output is a perfect matching, so
that nobody is left over; and the output is stable. Everything rests on a
handful of simple invariants, which are best isolated once and then used
repeatedly.

\subsection{The invariants}

\begin{observation}\label{obs:1}
Once a woman is engaged, she is engaged for the rest of the execution.
Her partner may change, but she never becomes free again.
\end{observation}

\begin{proof}
Inspect Algorithm~\ref{alg:gs}. An engaged woman only ever appears in
line~9, where she swaps one partner for another. There is no line in which a
woman becomes free.
\end{proof}

\begin{observation}\label{obs:2}
If an engaged woman's partner changes, the new partner is strictly better for
her than the old one. Consequently, the sequence of partners a woman has over
the course of the execution is strictly increasing in her preference order.
\end{observation}

\begin{proof}
A woman changes partner only in line~9, and the guard on line~8 is exactly the
condition that she strictly prefers the new proposer.
\end{proof}

\begin{observation}\label{obs:3}
The sequence of women to whom a fixed man proposes is strictly decreasing in
his preference order.
\end{observation}

\begin{proof}
By line~4 he always proposes to the best woman he has \emph{not yet} proposed
to, and once proposed to, a woman is struck off his list forever.
\end{proof}

\begin{figure}[ht]
\centering
\figplaceholder{50mm}{The two ratchets. On the left, a man's preference list
drawn as a vertical column with an arrow moving strictly downward over time
(Observation~\ref{obs:3}); on the right, a woman's preference list with an
arrow moving strictly upward (Observation~\ref{obs:2}). A caption such as
``men get worse off, women get better off, and neither can ever go back''
makes the point. This single picture explains termination, the $n^2$ bound,
and half of the stability proof.}
\caption{The monotonicity that makes the algorithm work.}
\label{fig:ratchets}
\end{figure}

\begin{intuition}
Observations~\ref{obs:2} and~\ref{obs:3} are the ratchets promised earlier.
The state of the algorithm can be summarised by ``how far down his list each
man has got'', and this vector only ever increases. That is why the process
cannot cycle the way Example~\ref{ex:cycle} did, and it immediately bounds the
number of steps. Everything else in this section is bookkeeping around these
two facts.
\end{intuition}

\subsection{Termination and running time}

\begin{lemma}[Free men always have somewhere to go]\label{lem:notstuck}
If a man is free at some point during the execution, then there is a woman to
whom he has not yet proposed. Hence line~4 is always well defined and the
algorithm never gets stuck.
\end{lemma}

\begin{proof}
Suppose for contradiction that $m$ is free and has already proposed to all $n$
women. By Observation~\ref{obs:1}, every woman who has ever received a
proposal is engaged --- she accepted the first offer she ever got, since a
free woman always accepts, and she has stayed engaged ever since. So all $n$
women are currently engaged.

Now count. The $n$ women are engaged to $n$ distinct men, so all $n$ men are
engaged. In particular $m$ is engaged, contradicting the assumption that he is
free.
\end{proof}

\begin{figure}[ht]
\centering
\figplaceholder{42mm}{The counting argument of Lemma~\ref{lem:notstuck}.
Left column: $n$ men with $m$ singled out and marked ``free''. Right column:
all $n$ women, each drawn with a solid engagement edge. The picture should
make it visually absurd that $n$ women are matched to $n$ distinct men while
one man is left over --- perhaps with a large ``$n$ edges, $n$ men, but $m$ is
not among them?'' annotation.}
\caption{A free man cannot have exhausted his list: that would mean all women
are engaged, hence all men are engaged.}
\label{fig:counting}
\end{figure}

\begin{lemma}[Termination and step count]\label{lem:steps}
The algorithm performs at most $n^2$ proposals and then terminates.
\end{lemma}

\begin{proof}
By Observation~\ref{obs:3}, each man proposes to each woman at most once, so
there are at most $n \cdot n = n^2$ proposals in total. Every iteration of the
while-loop performs exactly one proposal, so there are at most $n^2$
iterations. By Lemma~\ref{lem:notstuck} no iteration can fail to find a target
for the proposal, so the loop runs until its guard fails, i.e.\ until every
man is matched.
\end{proof}

\begin{lemma}[The output is a perfect matching]\label{lem:perfect}
When the algorithm terminates, every man and every woman is matched, and the
set of engaged pairs is a perfect matching.
\end{lemma}

\begin{proof}
At every moment the engagements form a partial matching: each man is engaged
to at most one woman and vice versa, directly by the update rules. The loop
terminates only when no man is free, i.e.\ when all $n$ men are engaged; since
they are engaged to distinct women and there are only $n$ women, all women are
engaged too.
\end{proof}

\begin{remark}[The $n^2$ bound is tight]
It is not hard to construct instances on which almost all $n^2$ proposals are
made; a classic family has every man with the same preference list and the
women's lists arranged so that each proposal displaces somebody. Since the
input itself has size $\Theta(n^2)$ --- there are $2n$ lists of length $n$ ---
the algorithm is optimal up to constants in the worst case, simply because
reading the input already costs that much.
\end{remark}

\subsection{Implementing it in $\bigoh{n^2}$ time}\label{sec:impl}

Lemma~\ref{lem:steps} bounds the number of \emph{proposals}, which is not
quite the same as bounding the running time. To get $\bigoh{n^2}$ actual
work, each proposal must be handled in $\bigoh{1}$ time. Three small data
structures suffice.

\begin{itemize}
  \item \textbf{A stack (or queue) of free men.} Maintaining ``the first
        unmatched man in the ordering'' by scanning would cost $\bigoh{n}$ per
        step. Instead keep the free men in any container with $\bigoh{1}$
        insert and extract. When a man is jilted, push him back. (Since the
        output is independent of the order, we may as well use whichever
        container is cheapest.)
  \item \textbf{A ``next proposal'' pointer per man.} Man $m$ stores an index
        into his own preference list, initially $1$, incremented after every
        proposal. By Observation~\ref{obs:3} this is exactly the information
        needed for line~4, and it makes that line $\bigoh{1}$.
  \item \textbf{A rank matrix for the women.} The comparison on line~8 asks
        ``does $w$ prefer $m$ to $m''$?'', which naively means searching her
        list: $\bigoh{n}$. Precompute a table $\mathrm{rank}[w][m]$ giving the
        position of $m$ on $w$'s list. Then the test is just
        $\mathrm{rank}[w][m] < \mathrm{rank}[w][m'']$, in $\bigoh{1}$.
        Building the table costs $\bigoh{n^2}$, which we are paying anyway to
        read the input.
\end{itemize}

With these, every proposal costs $\bigoh{1}$ and the total running time is
$\bigoh{n^2}$, matching the input size. The space usage is $\bigoh{n^2}$,
dominated by the preference lists and the rank matrix.

\subsection{Stability}

The remaining --- and most important --- claim is that the matching we produce
cannot be undermined.

\begin{theorem}[Correctness]\label{thm:stable}
The matching returned by Algorithm~\ref{alg:gs} is stable.
\end{theorem}

\begin{proof}
Let $\mu$ be the output. Consider any man $m$ and any woman $w'$ with
$\mu(m) = w \neq w'$; write $m' = \mu(w')$. We must show that $(m,w')$ is not
a blocking pair, i.e.\ that at least one of the two people involved is not
interested. We distinguish two cases according to whether $m$ ever proposed to
$w'$.

\smallskip
\textbf{Case 1: $m$ never proposed to $w'$.}
The algorithm terminated with $m$ engaged to $w$, and a man becomes engaged
only by proposing; so $m$ did propose to $w$ at some point. By
Observation~\ref{obs:3} the women he proposed to form a prefix of his
preference list, hence every woman he strictly prefers to $w$ received a
proposal from him. Since $w'$ did not, we conclude $w \pr{m} w'$: he prefers
his own partner to $w'$, and does not wish to deviate.

\smallskip
\textbf{Case 2: $m$ did propose to $w'$.}
Since $m$ is not matched to $w'$ at the end, one of two things happened: $w'$
rejected him on the spot, or she accepted and later left him for somebody
better. In the first case she was, at that moment, engaged to a man she
preferred to $m$; in the second case she immediately afterwards became engaged
to a man she preferred to $m$ (Observation~\ref{obs:2}). Either way, from that
moment on $w'$ was engaged to somebody she strictly prefers to $m$.

By Observation~\ref{obs:1} she never becomes free again, and by
Observation~\ref{obs:2} her partner only ever improves. Hence her final
partner $m' = \mu(w')$ satisfies $m' \pr{w'} m$: she prefers her own partner
to $m$, and does not wish to deviate.

\smallskip
In both cases $(m,w')$ fails to be a blocking pair. As $m$ and $w'$ were
arbitrary, $\mu$ is stable.
\end{proof}

\begin{figure}[ht]
\centering
\figplaceholder{55mm}{The stability proof in one picture. Draw the two matched
pairs $m$--$w$ (top) and $m'$--$w'$ (bottom), with the potential blocking edge
$m$--$w'$ dashed across the middle. Then two annotated branches: Case~1 ``$m$
never proposed to $w'$'' with the label ``$\Rightarrow w \succ_m w'$: he is
not interested'', and Case~2 ``$m$ proposed and was turned down'' with the
label ``$\Rightarrow m' \succ_{w'} m$: she is not interested''. This is the
figure students remember; it is worth making it large.}
\caption{Why no blocking pair can survive: either he never asked (so he does
not want her), or he asked and was refused (so she does not want him).}
\label{fig:stabilityproof}
\end{figure}

\begin{intuition}
The proof is a case distinction with exactly two cases, and each is one line
of reasoning. \emph{If he never asked her, it is because he is doing better
than her already} --- that is Observation~\ref{obs:3}, the men's ratchet.
\emph{If he did ask and is not with her, it is because she found somebody
better and never went back} --- that is Observations~\ref{obs:1}
and~\ref{obs:2}, the women's ratchet. Every blocking pair requires both people
to be interested, and the algorithm's asymmetry guarantees that at least one
of them is not.

It is also instructive to see where this argument would break for the
roommates problem of Section~\ref{sec:roommates}: with no two sides, there is
no coherent way to say who proposes and who disposes, so neither ratchet can
be set up. The proof knows about bipartiteness, as it must.
\end{intuition}

\begin{proof}[Proof of Theorem~\ref{thm:main}]
Lemma~\ref{lem:steps} gives termination after at most $n^2$ proposals,
Lemma~\ref{lem:perfect} says that the output is a perfect matching, and
Theorem~\ref{thm:stable} says that it is stable. Since the algorithm always
succeeds, a stable matching exists for every instance.
\end{proof}

%==============================================================================
\section{Who wins? Optimality of the proposing side}\label{sec:optimality}
%==============================================================================

We now know a stable matching always exists and can be found in $\bigoh{n^2}$
time. Two questions immediately suggest themselves, and they turn out to be
the same question:

\begin{itemize}[nosep]
  \item Can there be more than one stable matching? If so, which one do we get?
  \item The algorithm is blatantly asymmetric --- men act, women react. Does it
        matter who proposes?
\end{itemize}

\subsection{It does matter}

\begin{example}[The battle of the sexes]\label{ex:2x2}
Take $n = 2$:
\[
  \text{Arie}:\ \text{Ann} \succ \text{Betty}, \qquad
  \text{Bert}:\ \text{Betty} \succ \text{Ann},
\]
\[
  \text{Ann}:\ \text{Bert} \succ \text{Arie}, \qquad
  \text{Betty}:\ \text{Arie} \succ \text{Bert}.
\]
Each man's first choice is a woman whose first choice is the other man. The
preferences are as opposed as they can possibly be.
\end{example}

There are only two matchings here, and both are stable.

\begin{itemize}
  \item $\mu_{\Men} = \{(\text{Arie},\text{Ann}), (\text{Bert},\text{Betty})\}$:
        both men get their first choice, both women get their last. There is no
        blocking pair, because neither man wants to move.
  \item $\mu_{\Women} = \{(\text{Arie},\text{Betty}), (\text{Bert},\text{Ann})\}$:
        both women get their first choice, both men get their last. Again no
        blocking pair, because neither woman wants to move.
\end{itemize}

Run Algorithm~\ref{alg:gs} with the men proposing: Arie asks Ann, Bert asks
Betty, both are free and accept, done --- we get $\mu_{\Men}$. Run the mirror
image with the women proposing and we get $\mu_{\Women}$. Same instance, same
algorithm, opposite outcomes. Every participant's fate is decided entirely by
which side was given the right to propose.

\begin{figure}[ht]
\centering
\figplaceholder{45mm}{The two stable matchings of Example~\ref{ex:2x2}, side
by side. Left: $\mu_{\Men}$, with little happy faces on the men and sad faces
on the women; right: $\mu_{\Women}$, the reverse. Below, an axis labelled
``men's welfare'' pointing one way and ``women's welfare'' pointing the other,
to convey that the two are in direct opposition.}
\caption{The same instance, both stable matchings. Whoever proposes, wins.}
\label{fig:2x2}
\end{figure}

\begin{remark}
Compare this with Example~\ref{ex:main}, where the men-proposing run produced
$\mu^*$. If you run the women-proposing version on that instance (do it --- it
takes five proposals), you get $\mu^*$ again. As we shall see in
Corollary~\ref{cor:unique}, this means Example~\ref{ex:main} has exactly
\emph{one} stable matching, which is why the choice of proposer made no
difference there.
\end{remark}

So the general question ``is the outcome unique?'' has the answer ``no''. What
we can hope for instead is a precise description of \emph{which} stable
matching the algorithm produces. The answer is as clean as one could wish.

\subsection{Best and worst stable partners}

\begin{definition}[Stable partner]
A woman $w$ is a \emph{stable partner} of a man $m$ if there exists \emph{some}
stable matching $\mu$ with $\mu(m) = w$. Since at least one stable matching
exists (Theorem~\ref{thm:main}), every participant has at least one stable
partner.

Write $\mathrm{best}(m)$ for the woman highest on $m$'s preference list among
his stable partners, and $\mathrm{worst}(w)$ for the man lowest on $w$'s list
among her stable partners.
\end{definition}

These are definitions about the \emph{set of all stable matchings} --- a set
which may be exponentially large and which we certainly do not want to
enumerate. Note in particular what is \emph{not} being claimed: there is no
reason a priori why the assignment $m \mapsto \mathrm{best}(m)$ should even be
a matching. Two different men could conceivably have the same best stable
partner, in which case ``give everybody their best stable partner
simultaneously'' would be incoherent. The following theorem says that this
never happens, and, moreover, that Gale--Shapley computes it.

\begin{theorem}[Man-optimality]\label{thm:manoptimal}
Every execution of Algorithm~\ref{alg:gs} --- for every ordering of the men
and every choice of which free man proposes next --- returns the same set
\[
  \mu_{\Men} \;=\; \bigl\{\,(m,\ \mathrm{best}(m)) \;:\; m \in \Men \,\bigr\}.
\]
\end{theorem}

\begin{intuition}
Before the proof, here is the shape of the argument. A man ends up with the
first woman on his list who does not reject him, so the only way he can fail
to get his best stable partner is if some stable partner of his rejects him
along the way. We show that \emph{no stable partner ever rejects anybody}.

The proof is by ``consider the first time something goes wrong''. Suppose some
stable partner does reject somebody, and look at the very first such rejection
in the execution. Being first is powerful: it means that up to this moment the
algorithm has made no mistakes at all, so anything we know about the state of
the algorithm can be leveraged. What we extract is enough information to build
a blocking pair inside a stable matching --- a contradiction.
\end{intuition}

\begin{proof}
By Theorem~\ref{thm:stable} the output $\mu$ of any execution is a stable
matching, so $\mu(m)$ is always \emph{a} stable partner of $m$. It therefore
suffices to prove that $\mu(m)$ is at least as good as $\mathrm{best}(m)$ for
every $m$; then $\mu(m) = \mathrm{best}(m)$ and, as this holds for every
execution, all executions agree.

\medskip
\noindent\textbf{Main claim: during the execution, no man is ever rejected by
a stable partner.}

Here ``rejected by $w$'' covers both possibilities: $w$ turning down a
proposal, and $w$ breaking off an engagement in favour of a better proposer.

Suppose the claim fails. Fix the \emph{first} moment in the execution at which
a man is rejected by a stable partner of his, say
\[
  \text{$w$ rejects $m$, and $w$ is a stable partner of $m$.}
\]
Because $w$ is a stable partner of $m$, there is a stable matching $S$ with
$S(m) = w$. Because $w$ rejects $m$, she does so in favour of some man $m'$
whom she strictly prefers:
\begin{equation}\label{eq:opt1}
  m' \pr{w} m .
\end{equation}
(If she rejected a proposal, $m'$ is her current partner; if she broke an
engagement, $m'$ is the new proposer.)

Let $w' = S(m')$ be the partner of $m'$ in the stable matching $S$. Note
$w' \neq w$, since $S(w) = m \neq m'$.

\medskip
\noindent\emph{Sub-claim: $w \pr{m'} w'$.}

At the moment under consideration $m'$ has just proposed to $w$ (or is engaged
to her, having proposed earlier). By Observation~\ref{obs:3}, every woman whom
$m'$ strictly prefers to $w$ has already rejected him at some earlier moment.
Now suppose, for contradiction, that $w' \pr{m'} w$. Then $w'$ is one of those
women: she has already rejected $m'$ earlier in the execution. But $w'$ is a
stable partner of $m'$, witnessed by $S$. So that earlier event was a rejection
of a man by a stable partner --- contradicting the choice of our moment as the
\emph{first} such event. Hence $w \pr{m'} w'$, proving the sub-claim.

\medskip
Now put the pieces together in the stable matching $S$, where $m'$ is matched
to $w'$ and $w$ is matched to $m$:
\begin{itemize}[nosep]
  \item $w \pr{m'} w'$ \quad (sub-claim): $m'$ prefers $w$ to his $S$-partner;
  \item $m' \pr{w} m$ \quad \eqref{eq:opt1}: $w$ prefers $m'$ to her
        $S$-partner.
\end{itemize}
So $(m', w)$ is a blocking pair for $S$, contradicting the stability of $S$.
This proves the main claim.

\medskip
\noindent\textbf{From the claim to the theorem.}
Consider the output $\mu$ and any man $m$. He is matched to the first woman on
his list who did not reject him: by Observation~\ref{obs:3} he worked strictly
downward from the top, and every woman above $\mu(m)$ on his list rejected him
at some point. By the main claim, none of those women is a stable partner of
$m$. Hence $\mathrm{best}(m)$ is not among them, so $\mathrm{best}(m)$ is at
position $\mu(m)$ or below on $m$'s list; that is, $\mu(m)$ is at least as
good for $m$ as $\mathrm{best}(m)$. Since $\mu(m)$ \emph{is} a stable partner
of $m$ and $\mathrm{best}(m)$ is by definition his best one, we get
$\mu(m) = \mathrm{best}(m)$.
\end{proof}

\begin{figure}[ht]
\centering
\figplaceholder{60mm}{The man-optimality proof. Suggested layout: a horizontal
time axis for the execution, with a marked point ``first rejection of a man by
a stable partner: $w$ rejects $m$ in favour of $m'$''. Everything to the left
shaded as ``no stable partner has rejected anyone yet''. Below the axis, draw
the stable matching $S$ as two solid edges $m$--$w$ and $m'$--$w'$, and mark
the blocking pair $(m',w)$ as a red dashed edge, annotated with the two
preferences $w \succ_{m'} w'$ and $m' \succ_w m$. Arrows connecting the two
halves would show how the timeline fact yields the preference on the left.}
\caption{The first rejection by a stable partner would create a blocking pair
in a stable matching.}
\label{fig:optproof}
\end{figure}

\begin{corollary}\label{cor:1}
The set $\{(m, \mathrm{best}(m)) : m \in \Men\}$ is a matching, and it is
stable.
\end{corollary}

\begin{proof}
It is the output of Algorithm~\ref{alg:gs}, which is a matching by
Lemma~\ref{lem:perfect} and stable by Theorem~\ref{thm:stable}.
\end{proof}

This is worth savouring. Each man selfishly names the best partner he could
have in any stable matching whatsoever --- ignoring everybody else's
constraints entirely --- and these individually optimal, apparently
uncoordinated demands turn out to be simultaneously satisfiable, and to form a
stable matching. There is no reason to expect this from the definitions; it
falls out of the algorithm. This is a good example of a theorem that is
\emph{much} easier to prove with an algorithm than without one.

\begin{corollary}\label{cor:2}
Algorithm~\ref{alg:gs} produces a \emph{man-optimal} stable matching: no
stable matching gives any man a strictly better partner. In particular its
output does not depend on the order in which the men propose.
\end{corollary}

\subsection{And now the bad news}

Man-optimality sounds like an unambiguous success story. It is not, because
the gain of one side is exactly the loss of the other.

\begin{corollary}[Woman-pessimality]\label{cor:3}
The matching produced by Algorithm~\ref{alg:gs} is the \emph{worst} stable
matching from the women's point of view:
\[
  \mu_{\Men} = \bigl\{\,(\mathrm{worst}(w),\ w) \;:\; w \in \Women \,\bigr\}.
\]
That is, every woman simultaneously receives the worst partner she has in any
stable matching.
\end{corollary}

\begin{proof}
Write $\mu = \mu_{\Men}$ and suppose the statement fails: there is a woman $w$
with $\mu(w) = m$ such that $m$ is \emph{not} her worst stable partner. Then
some stable matching $S$ gives her a strictly worse partner
$m' = S(w)$, so
\begin{equation}\label{eq:cor3a}
  m \pr{w} m' .
\end{equation}
Let $w' = S(m)$ be the partner of $m$ in $S$. Since $S(w) = m' \neq m$, we have
$w' \neq w$.

Now, $w'$ is a stable partner of $m$ (witnessed by $S$), and by
Theorem~\ref{thm:manoptimal} the woman $w = \mu(m)$ is his \emph{best} stable
partner. As $w' \neq w$, this gives
\begin{equation}\label{eq:cor3b}
  w \pr{m} w' .
\end{equation}
But now look at $S$, in which $m$ is matched to $w'$ and $w$ is matched to
$m'$. By~\eqref{eq:cor3b} $m$ prefers $w$ to his $S$-partner, and
by~\eqref{eq:cor3a} $w$ prefers $m$ to her $S$-partner. So $(m,w)$ is a
blocking pair for $S$, contradicting the stability of $S$.
\end{proof}

\begin{figure}[ht]
\centering
\figplaceholder{50mm}{The proof of Corollary~\ref{cor:3}. Two panels. Left
panel, the Gale--Shapley output $\mu$: edge $m$--$w$, annotated ``$w$ is $m$'s
\emph{best} stable partner (Theorem~\ref{thm:manoptimal})''. Right panel, the
hypothetical stable matching $S$: edges $m$--$w'$ and $m'$--$w$, with the red
dashed blocking edge $m$--$w$ across it and the two labels $w \succ_m w'$ and
$m \succ_w m'$. An arrow from the left panel to the right showing where
$w \succ_m w'$ comes from.}
\caption{If some woman could do worse in a stable matching than
Gale--Shapley gives her, that matching would not be stable.}
\label{fig:cor3proof}
\end{figure}

\begin{intuition}
The mechanism behind this is easier to feel than to state. The men-proposing
algorithm has the men working \emph{down} their lists, stopping as early as
possible; it has the women working \emph{up} their lists, but only among the
men who happen to knock on their door. A woman can never reach out to a man
she prefers; she can only wait. So she ends up with the best of a set of
offers that was itself curated by the men's optimisation. Symmetrically, in
the women-proposing algorithm every woman gets her best stable partner and
every man his worst.

More generally, the interests of the two sides are exactly opposed across the
set of stable matchings: whenever one stable matching is better than another
for some man, it is worse for the woman involved. (The set of all stable
matchings in fact forms a distributive lattice, with $\mu_{\Men}$ and
$\mu_{\Women}$ as its two extreme elements. That structure is beyond this
lecture, but it is the right mental picture: a whole ordered family of
compromises stretching between the two algorithmic extremes.)
\end{intuition}

\begin{figure}[ht]
\centering
\figplaceholder{45mm}{Optional but useful mental picture: the set of all
stable matchings drawn as a lattice diagram (a diamond-ish Hasse diagram),
with $\mu_{\Men}$ at the top labelled ``best for all men / worst for all
women'' and $\mu_{\Women}$ at the bottom labelled ``worst for all men / best
for all women'', and a few unnamed compromise matchings in between. Mark
clearly that only the two extremes are computed by the algorithms of this
lecture.}
\caption{Between the two extremes computed by Gale--Shapley there may be many
other stable matchings.}
\label{fig:lattice}
\end{figure}

\begin{corollary}[A uniqueness test]\label{cor:unique}
If the men-proposing and women-proposing executions return the same matching,
then it is the only stable matching of the instance.
\end{corollary}

\begin{proof}
By Theorem~\ref{thm:manoptimal} and its mirror image, $\mu_{\Men}$ gives every
man his best stable partner and $\mu_{\Women}$ gives every man his worst (by
Corollary~\ref{cor:3} applied with the roles of the sides exchanged). If
$\mu_{\Men} = \mu_{\Women}$, then for each man his best and worst stable
partners coincide, so he has exactly one stable partner, and the stable
matching is unique.
\end{proof}

This explains the observation made after Example~\ref{ex:2x2}: in
Example~\ref{ex:main} both executions return $\mu^*$, so $\mu^*$ is the unique
stable matching of that instance --- which is why the choice of proposer was
invisible there, and why running the algorithm with a different ordering of
the men changed nothing.

\subsection{What this means in practice}

The theory has an immediate design consequence, and it is not a subtle one:
\emph{choosing who proposes is choosing who wins}. A clearinghouse is not a
neutral computational device; it makes a distributional decision, and the
mathematics makes that decision unusually explicit.

This was not merely theoretical for the residency match. The algorithm the
NRMP had used since 1952 was, in the terms of this lecture, hospital-proposing:
it produced the hospital-optimal, and hence applicant-pessimal, stable
matching. Once Roth's analysis made this visible, the direction of the
algorithm became a subject of open dispute, and in the 1998 redesign by Roth
and Peranson the match was switched to applicant-proposing. (The real system
also has to handle couples who want positions in the same city, programmes
with many positions rather than one, and applicants who rank only part of the
list --- so it is not literally Algorithm~\ref{alg:gs}. But the core is
deferred acceptance, and the choice of which side proposes is exactly the
choice analysed above.)

\begin{figure}[ht]
\centering
\figplaceholder{40mm}{Two identical clearinghouse boxes, one labelled
``hospitals propose'' with an arrow to a happy hospital and a resigned
applicant, the other labelled ``applicants propose'' with the reverse. Caption
underneath: ``the same data, the same algorithm, a different answer to who
gets the benefit of the doubt''.}
\caption{The direction of proposal is a policy decision, not a technical
detail.}
\label{fig:whoproposes}
\end{figure}

A second practical question follows immediately, and we only pose it here: if
the mechanism systematically disadvantages the receiving side, do participants
have an incentive to \emph{lie} about their preferences? The answer, in brief,
is that the proposing side never benefits from lying, but the receiving side
sometimes does. Exercise~\ref{exc:lying} works out the smallest example.

%==============================================================================
\section{Summary and outlook}\label{sec:outlook}
%==============================================================================

\subsection*{What we proved}

\begin{itemize}
  \item Every stable marriage instance has a stable matching, and one can be
        found in $\bigoh{n^2}$ time by the Gale--Shapley algorithm
        (Theorem~\ref{thm:main}). Since the input has size $\Theta(n^2)$, this
        is optimal.
  \item The algorithm is highly non-deterministic --- any free man may propose
        next --- yet its output is completely determined by the instance
        (Theorem~\ref{thm:manoptimal}).
  \item That output is extreme in both directions at once: every man gets the
        best partner he has in any stable matching, and every woman
        simultaneously gets the worst (Corollary~\ref{cor:2} and
        Corollary~\ref{cor:3}).
  \item Existence is genuinely two-sided: for the stable roommates problem,
        where the bipartite structure is absent, stable matchings need not
        exist at all (Example~\ref{ex:roommates}).
\end{itemize}

\subsection*{Where the subject goes next (non-examinable)}

Stable matching turned out to be one of those rare problems that stays
tractable under a remarkable amount of generalisation, while a few natural
variants fall off a cliff into NP-hardness. A sampler:

\begin{itemize}
  \item \textbf{Unequal sides and incomplete lists.} If participants may
        declare some partners unacceptable, stable matchings still exist and
        are still found by essentially the same algorithm, but need no longer
        be perfect. The \emph{rural hospitals theorem} says something striking:
        the set of participants who go unmatched is the same in \emph{every}
        stable matching. Unpopular rural hospitals cannot fix their staffing
        problem by lobbying for a different stable outcome.
  \item \textbf{Many-to-one (hospitals/residents).} Hospitals have capacities.
        Everything in this lecture goes through with the obvious modifications;
        this is the version actually deployed in residency matching and in
        school choice systems.
  \item \textbf{Ties (indifference).} If preference lists may contain ties, one
        must distinguish weak, strong and super-stability. Some of the
        resulting problems remain polynomial; finding a \emph{maximum} weakly
        stable matching becomes NP-hard.
  \item \textbf{Forbidden pairs, group constraints, couples.} Many such
        variants remain solvable; the couples version used by the real NRMP
        can fail to admit a stable matching, and is NP-hard in general.
  \item \textbf{Strategy.} Men-proposing Gale--Shapley is
        \emph{strategy-proof for the men}: no man can ever gain by
        misreporting his list. No mechanism can be strategy-proof for both
        sides at once, and the receiving side can indeed gain by lying
        (Exercise~\ref{exc:lying}).
  \item \textbf{Stable roommates.} Despite Example~\ref{ex:roommates}, Irving
        gave an $\bigoh{n^2}$ algorithm that finds a stable matching or
        correctly reports that none exists. Adding indifference to the
        roommates problem makes deciding existence NP-complete.
  \item \textbf{How many stable matchings can there be?} Exponentially many:
        instances with $\Omega(2.28^n)$ stable matchings are known, building on
        a construction of Irving and Leather. An upper bound of the form $c^n$
        for an absolute constant $c$ was open for decades and was finally
        proved by Karlin, Oveis Gharan and Weber in 2018
        \cite{KarlinOveisGharanWeber2018}.
        %% NOTE TO SELF: slide 17 states a specific constant that came out of
        %% the PDF text extraction as "(3.55)". Check the original slide and
        %% put the intended constant here before distributing the notes.
\end{itemize}

%==============================================================================
\section{Exercises}
%==============================================================================

\begin{exercise}
Verify the remaining three steps of the cycle in Example~\ref{ex:cycle}: check
in each case that the indicated pair really is blocking, and that resolving it
gives the stated matching. Then find the (unique?) stable matching of that
instance by running Gale--Shapley.
\end{exercise}

\begin{exercise}
Run the women-proposing version of Gale--Shapley on Example~\ref{ex:main} and
confirm that it returns $\mu^*$. What does Corollary~\ref{cor:unique} let you
conclude?
\end{exercise}

\begin{exercise}
Give an instance with $n = 4$ that has at least four stable matchings.
\emph{Hint:} take two disjoint copies of Example~\ref{ex:2x2} and make each
participant prefer everybody in their own copy to everybody in the other.
Generalise to show that there are instances with at least $2^{n/2}$ stable
matchings.
\end{exercise}

\begin{exercise}
Construct a family of instances on which the men-proposing algorithm makes
$\Omega(n^2)$ proposals, showing that the bound of Lemma~\ref{lem:steps} is
tight up to a constant factor.
\end{exercise}

\begin{exercise}
Where exactly does the proof of Theorem~\ref{thm:stable} use that the
preference lists are strict (no ties)? Where does it use that they are
complete?
\end{exercise}

\begin{exercise}
Prove directly from the definitions --- without invoking
Theorem~\ref{thm:manoptimal} --- that if $\mu$ and $\nu$ are stable matchings
and some man strictly prefers $\mu$ to $\nu$, then his partner in $\mu$
strictly prefers $\nu$ to $\mu$. (This is the ``opposed interests'' statement
that underlies the lattice structure.)
\end{exercise}

\begin{exercise}\label{exc:lying}
Consider Example~\ref{ex:2x2}, and allow participants to submit lists that
declare some partners unacceptable. Suppose Ann submits the list ``Bert''
only, and everybody else reports truthfully. Run the men-proposing algorithm
on the reported preferences and check that Ann ends up with Bert --- her true
first choice, and strictly better than what truthful reporting gave her.
Explain in one sentence why this trick is available to the receiving side but
not to the proposing side.
\end{exercise}

\begin{exercise}
Implement Gale--Shapley as described in Section~\ref{sec:impl} and verify
empirically, on random instances with $n = 1000$, roughly how many proposals
are actually made and how good the men's and women's average partner ranks
are. (The behaviour on random instances is strikingly asymmetric; the men do
about $\ln n$ times better than the women.)
\end{exercise}

%==============================================================================
\begin{thebibliography}{9}

\bibitem{GaleShapley1962}
D.~Gale and L.~S.~Shapley.
\newblock College admissions and the stability of marriage.
\newblock \emph{The American Mathematical Monthly}, 69(1):9--15, 1962.

\bibitem{Knuth1976}
D.~E.~Knuth.
\newblock \emph{Mariages stables et leurs relations avec d'autres problèmes
combinatoires}.
\newblock Les Presses de l'Université de Montréal, 1976.

\bibitem{GusfieldIrving1989}
D.~Gusfield and R.~W.~Irving.
\newblock \emph{The Stable Marriage Problem: Structure and Algorithms}.
\newblock MIT Press, 1989.

\bibitem{Roth1984}
A.~E.~Roth.
\newblock The evolution of the labor market for medical interns and residents:
a case study in game theory.
\newblock \emph{Journal of Political Economy}, 92(6):991--1016, 1984.

\bibitem{RothSotomayor1990}
A.~E.~Roth and M.~Sotomayor.
\newblock \emph{Two-Sided Matching: A Study in Game-Theoretic Modeling and
Analysis}.
\newblock Cambridge University Press, 1990.

\bibitem{KarlinOveisGharanWeber2018}
A.~R.~Karlin, S.~Oveis Gharan and R.~Weber.
\newblock A simply exponential upper bound on the maximum number of stable
matchings.
\newblock In \emph{Proceedings of STOC 2018}, pages 920--925.

\end{thebibliography}

\end{document}
